\section{Introduction}
\label{sec:intro}

Reconstructing 3D objects from monocular images is a fundamental problem in computer vision. An efficient reconstruction system opens up a wide range of applications, including augmented reality, filmmaking, and manufacturing. Monocular 3D reconstruction is also a complex inverse problem: while the visible surface can be estimated from shading, predicting the occluded surface necessitates a strong 3D object prior. Our field has seen a divergence in two different directions: feedforward regression~\cite{choy20163d,wang2018pixel2mesh,mescheder2019occupancy,xu2019disn,groueix2018,wu2017marrnet,huang2023shapeclipper,wu2023multiview,huang2024zeroshape,hong2023lrm,TripoSR2024,tang2024lgm,zhang2024gs,xu2024instantmesh,wang2024crm,sf3d2024} and diffusion-based generation~\cite{nichol2022point,jun2023shap,liu2023one2345++,shue20233d,chou2023diffusion,shim2023diffusion,cheng2023sdfusion,li2023diffusion,liu2023zero,shi2023mvdream,liu2023syncdreamer,huang2024pointinfinity,yariv2024mosaic,chen2023single,zhao2023michelangelo,lan2024ln3diff}. Despite the significant progress made in both directions, each has fundamental limitations.

Regression-based models are highly effective in adhering to the visible surface in the image, and the inference speed is typically fast. However, they make the oversimplified assumption of bijective mapping between images and 3D. This assumption introduces ambiguity in the learning objective, leading to poorly estimated surfaces and textures in occluded regions. On the other hand, diffusion-based approaches are generative and do not predict the statistical mean. However, their iterative sampling at inference time is computationally inefficient when modeling high-resolution 3D. Additionally, previous studies such as~\cite{huang2024zeroshape} indicate that diffusion-generated 3D models exhibit worse alignment to the surface visible in the input image. How can we take the best of both worlds while avoiding their limitations?

In light of this, we propose \method, which breaks the 3D reconstruction process down into two stages: the point sampling stage and the meshing stage. The point sampling stage uses diffusion models to generate sparse point clouds, followed by the meshing stage, transforming point clouds into highly detailed meshes. Our main idea is to offload the uncertainty modeling to the point sampling stage, where the low resolution of the point clouds allows rapid iterative sampling. The subsequent meshing stage leverages the local image features to transform the point cloud into a detailed mesh of high output fidelity. Reducing the meshing uncertainty with point clouds further facilitates unsupervised learning of inverse rendering, which reduces the baked-in lighting in the textures. Our two-stage design enables \method to significantly outperform previous regressive methods, while preserving high computational efficiency and fidelity to input observation.

A key design choice of our method is the usage of point clouds to connect the two stages. To ensure fast reconstruction, our intermediate representation needs to be lightweight so it can be efficiently generated. However, it should provide enough guidance to the meshing stage. This inspires us to use point clouds, which are perhaps the most computationally efficient 3D representation because all information bits are used to represent the surface. Moreover, the lack of connectivity, typically considered as the drawback of point clouds, now turns into an advantage with our two-stage approach for editing purposes. When the back surface does not align with user expectations, local edits can be easily made on the low-resolution point clouds without worrying about topologies (see~\cref{fig:teaser} bottom). Feeding edited point clouds into the meshing stage produces better meshes tailored towards user requirements.

Our experiments demonstrate the superiority of \method over previous state-of-the-art methods, with solid quantitative and qualitative results on various data sources.
\method also exhibits a strong generalization ability to in-the-wild images and AI-generated images. With a total inference time below \gentime seconds, \method is not only efficient but also allows for easy user-driven edits, offering a practical solution to the task of monocular 3D reconstruction. We hope that this is a meaningful step towards scalable generation of high-quality 3D assets.
